% Procedencia: integral 2249, §25.10--13 y cap. 89; propietarios 02_pantalla_accion,
% 37_cuerdas_branas_teoria_m, terna_espectral_teoria_m, capitulo_25_pantallas_teoria_m.
\section{Action on trajectories and eleven-dimensional dynamical realization}
\label{iv:accion-realizacion}

Extending the structure to dimensional dynamics requires
joint transport of the state, channels and boundary
data. This section begins with a functional defined on
TPK trajectories and its exact composition law.
The worldsheet and eleven-dimensional codomain are
introduced afterwards, with the maps that must preserve that
composition and the operators already constructed.

\subsection{Descent of an evolution and sufficiency of the observed datum}

Let \(q:\mathcal X_{\mathrm{HMT}}\to Y\) be a surjective
readout of states with memory, and let
\(U:\mathcal X_{\mathrm{HMT}}\to\mathcal X_{\mathrm{HMT}}\)
be a complete evolution. An evolution on observed
values exists precisely when the datum preserved by
\(q\) suffices to determine the next observed value.

\begin{theorem}[Exact criterion for descent along fibres]
\label{iv:thm:descenso-fibras}
There exists a unique map \(\overline U:Y\to Y\)
satisfying \(qU=\overline Uq\) if and only if
\begin{equation}
 q(x)=q(y)\quad\Longrightarrow\quad q(Ux)=q(Uy).
 \label{iv:eq:descenso-fibras}
\end{equation}
\end{theorem}
\begin{proof}
If \(\overline U\) exists, applying the same function to
\(q(x)=q(y)\) proves the condition. Conversely,
define \(\overline U(q(x))=q(Ux)\). Condition
\eqref{iv:eq:descenso-fibras} guarantees independence of the
representative. Surjectivity of \(q\) defines the
map on all of \(Y\) and proves uniqueness.
\end{proof}

The theorem determines what memory a readout
must preserve. When two representatives have different observed
images, the retained datum may be enlarged, the domain
restricted, or a multivalued correspondence used.
These are distinct mathematical choices. In particular,
return of a phase coordinate does not authorize identification
of the evolutions of two histories retaining different
memories.

\subsection{Cocycle of directional and arithmetic changes}

An admissible TPK trajectory carries, on each edge
\(e_j\), a direction \(d_j\) and an arithmetic type
\(t_j\). To count changes in a concatenation,
the preceding datum is also preserved. The augmented edge is
\(\widetilde e_j=(e_j,d_{j-1},t_{j-1})\); the pair
\((d_0,t_0)\) belongs to the initial boundary. Define
\[
 w_{\mathrm{giro}}(\widetilde e_j)
       =\mathbf1_{\{d_j\ne d_{j-1}\}},\qquad
 w_{\mathrm{tipo}}(\widetilde e_j)
       =\mathbf1_{\{t_j\ne t_{j-1}\}}.
\]
For \(\widetilde\gamma=(\widetilde e_1,\ldots,
\widetilde e_R)\), let
\begin{equation}
 \mathbf I(\widetilde\gamma)=
 \left(\sum_{j=1}^R w_{\mathrm{giro}}(\widetilde e_j),
       \sum_{j=1}^R w_{\mathrm{tipo}}(\widetilde e_j)\right)
 =\bigl(N_{\mathrm{giro}},N_{\mathrm{tipo}}\bigr).
 \label{iv:eq:cociclo-accion}
\end{equation}
Two trajectories compose when their endpoints agree
and the preceding datum of the second agrees with the terminal
datum of the first.

\begin{proposition}[Additivity with preserved boundary]
\label{iv:prop:accion-aditiva}
For every composable pair,
\begin{equation}
 \mathbf I(\widetilde\gamma\widetilde\eta)
 =\mathbf I(\widetilde\gamma)+\mathbf I(\widetilde\eta).
 \label{iv:eq:accion-aditiva}
\end{equation}
Thus, \(\mathbf I\) is an additive cocycle with values
in \(\mathbb N^2\) on the category of augmented
trajectories. For every \(\lambda>0\), the functional
\[
 I_\lambda(\widetilde\gamma)
 =N_{\mathrm{giro}}(\widetilde\gamma)
      +\lambda N_{\mathrm{tipo}}(\widetilde\gamma)
\]
is also additive.
\end{proposition}
\begin{proof}
Concatenation does not change the preceding datum of any
edge: for the first edge of \(\widetilde\eta\),
the composition condition makes it agree with the terminal datum
of \(\widetilde\gamma\). The sum over all edges
splits exactly into the two sums in
\eqref{iv:eq:accion-aditiva}. Applying the linear functional
\((a,b)\mapsto a+\lambda b\) proves the last statement.
\end{proof}

Boundary information is an effective part of the proof.
If two one-edge segments are evaluated separately
while ignoring the connection between them, both may have zero
turn count even though their directions differ. The
concatenated path then contains a turn. Preserving
\((d_{j-1},t_{j-1})\) records that term on the first
edge of the second segment and restores additivity.

For fixed endpoints and boundary data, the
minimum-interaction principle selects minimizers of
\(I_\lambda\) among admissible routes. The principle
is a selection rule, whereas the following
result establishes a finite-existence property.

\begin{proposition}[Finite existence of minimizers]
If the set \(\Gamma(a,b)\) of admissible trajectories
with the fixed boundary is finite and nonempty, there exists
\(\gamma_*\in\Gamma(a,b)\) such that
\(I_\lambda(\gamma_*)\leq I_\lambda(\gamma)\)
for all \(\gamma\in\Gamma(a,b)\).
\end{proposition}
\begin{proof}
The image of \(\Gamma(a,b)\) under \(I_\lambda\) is
a nonempty finite set of real numbers, which has
a minimum element and at least one preimage.
\end{proof}

The proposition does not select the value of \(\lambda\),
nor does it confuse existence with uniqueness. The value and
selection rule used in a particular realization
must be part of its declared data. The additive
functional recovered here quantifies changes of the path
itself; its transport to a dimensional action is
a later operation.

\subsection{Extended persistence and worldsheet}

Realization of a route as an extended object distributes
its genealogy along an internal coordinate. For a string,
the history has two local coordinates \(\sigma^a\):
one internal spatial coordinate and one temporal coordinate. A realization is
described by a map \(X:\Sigma\to\mathcal M\), a
metric \(h_{ab}\) on the worldsheet and background
fields in the target space. Conditions on
\(\partial\Sigma\) belong to the variational domain.
For a brane of spatial dimension \(p\), the history
is a worldvolume of dimension \(p+1\).

The variational realization of the string is expressed by
the Polyakov action
\begin{equation}
 S_{\mathrm P}[X,h]
 =-\frac{T_{\mathrm s}}2\int_\Sigma d^2\sigma\,
 \sqrt{-h}\,h^{ab}g_{MN}(X)
             \partial_aX^M\partial_bX^N,
 \label{iv:eq:polyakov}
\end{equation}
where \(T_{\mathrm s}\) is tension. The notation
\(T_{\mathrm s}\) distinguishes it from duality
\(\mathsf T_0\), cardinal transport \(T_d\)
and complete evolution \(U\).

Composition of augmented trajectories already has the
property proved in \ref{iv:prop:accion-aditiva}.
A continuous realization of that composition must
specify the map from paths to configurations,
its transport of boundary conditions and the
normalization comparing \(I_\lambda\) with the continuous
action. Integral \eqref{iv:eq:polyakov} specifies
the codomain and its fields; its parameters do not
retrospectively select TPK emissions or registers.

\subsection{Eleven-dimensional realization and graded operators}

The reduction directed by \(K\) provides
\(V_{11}=\ell_K\oplus V_{10}\), and Theorem
\ref{iv:thm:pantallas-isomorfas} coordinates this reduction
with the lattice screen and duality. The physical
realization is typed through
\[
 \Xi_M:\mathscr S_M^{\mathrm{alg}}
                  \longrightarrow\mathcal X_{11}^{\mathrm{phys}}.
\]
The codomain carries a metric \(g_{MN}\), a three-form
\(C_{MNP}\), a gravitino \(\psi_M\), an action,
gauge constraints and supersymmetry
transformations. An operator realization includes
\[
 \mathfrak S_M=(\mathcal H_{\mathrm{phys}},H,Q),\qquad
 \mathcal H_{\mathrm{phys}}=\mathcal H_+\oplus\mathcal H_-.
\]
The grading is expressed through an involution
\(\Gamma\), equal to \(+I\) on \(\mathcal H_+\)
and \(-I\) on \(\mathcal H_-\). The supercharge
is odd, \(\Gamma Q=-Q\Gamma\), and the
Hamiltonian is nonnegative. The relation
\(Q^2=H\) is understood as an equality of operators,
including their domains.

Let \(\mathcal H_{\mathrm{HMT}}\) be a Hilbert-space realization
of the relevant genealogical states,
with operators \(\widehat H,\widehat Q\). An
isometry \(W:\mathcal H_{\mathrm{HMT}}\to
\mathcal H_{\mathrm{phys}}\) must transport
domains, grading and observables. The operator
conditions are written
\begin{equation}
 W\widehat H=HW,\qquad
 W\widehat Q=QW,\qquad
 WU_T^{\mathrm{HMT}}=U_T^{\mathrm{phys}}W.
 \label{iv:eq:entrelazamiento-fisico}
\end{equation}

\begin{theorem}[Compatible transport of the supercharge]
\label{iv:thm:supercarga-transporte}
Suppose \(\widehat Q^{\,2}=\widehat H\),
\(W D(\widehat Q)\subset D(Q)\), and the
intertwining relations of \eqref{iv:eq:entrelazamiento-fisico}
hold on their domains. Then, for
\(\psi\in D(\widehat H)=D(\widehat Q^{\,2})\),
\[
 W\psi\in D(Q^2),\qquad Q^2W\psi=HW\psi.
\]
If physical transport of the radii also respects
\eqref{iv:eq:conjugacion-h}, the transported compact
sector preserves its spectrum under reciprocal radii.
\end{theorem}
\begin{proof}
Since \(\psi\in D(\widehat Q^2)\), both
\(\psi\) and \(\widehat Q\psi\) belong
to \(D(\widehat Q)\). Thus, \(W\psi\in D(Q)\)
and \(QW\psi=W\widehat Q\psi\in D(Q)\). One may
apply \(Q\) again, and
\[
 Q^2W\psi=W\widehat Q^{\,2}\psi
          =W\widehat H\psi=HW\psi.
\]
Spectral equality of the compact sector is transported
by the isometry and the unitary T-intertwiner;
on the closed image of \(W\), this is a
unitary equivalence of the corresponding realizations.
\end{proof}

The theorem explicitly preserves the operators and the
domain on which the identity is proved. The
grading organizes two sectors; construction of
fermionic fields and their occupation relations has
its own structure. Thermal distributions and the
species catalogue are not used in the proof
of the screens or in the compact duality of this
article.

\subsection{Compatible compactification and conclusion of the development}

A compactification of the state is described through
a prolongable section \(x_\infty\), a map
\(\kappa_{\mathrm{comp}}\) to the chosen moduli
space and a transport \(\mathcal U\) of the
metric, orientation, boundary and memory data. The
subscript distinguishes this map from the rational readout
\(\kappa\) of the dodecaphasic register. Its compatibility
requires four concrete properties:
\begin{enumerate}
\item the finite restrictions of \(x_\infty\)
      belong to admissible TPK prolongations;
\item the compactification map intertwines the dualities,
      \(\kappa_{\mathrm{comp}}\mathsf T
        =\mathsf T_M\kappa_{\mathrm{comp}}\);
\item observables and the spectrum descend to the
      chosen quotients according to criterion
      \ref{iv:thm:descenso-fibras};
\item the realization transports the direction
      \(\ell_K\), boundary conditions and operator
      domains, and preserves the declared action
      and supersymmetry equations.
\end{enumerate}

The preceding proofs provide the phase Hilbert-space
realization, tracial closure, screens
with their quotient maps, radius involution,
its unitary equivalence and the additive trajectory
cocycle. The eleven-dimensional realization
specifies the fields and compatibilities it
uses when receiving that structure.

\paragraph{Conditions for physical identification.}
\label{iv:condiciones-realizacion-fisica}
Complete physical identification must additionally
specify the tension and action normalization,
construct the oscillator sector and its amplitudes,
control anomalies and establish a consistent
low-energy gravitational limit for each stated
compactification.
The string action and change of scale used here
are expressed in \eqref{iv:eq:polyakov}
and \eqref{iv:eq:normalizacion-cuerda}, respectively.
Theorem~\ref{iv:thm:supercarga-transporte}
proves the supercharge identity on the image of
\(W\), under the domains and intertwining relations
in its statement. The conditions for complete physical
identification do not alter this identity or the other
algebraic proofs: they delimit the physical realization
to which those results are transported.

The exceptional and dimensional prolongations
remain linked by their common origin. The
Monster acts on the graded construction
\(V^\natural\); the compact Hamiltonian and
screens receive data from other maps of the
same state. This formulation preserves the
positive content of the dual realization,
with its operations, proofs and the information
transmitted at each stage: Theorem~\ref{exc:moonshine}
establishes the exceptional identification, and Theorems
\ref{iv:thm:t-unitaria} and~\ref{iv:thm:pantallas-isomorfas}
establish compact equivalence and screen coordination.
